\Problem{Bicolor numbers}
\Statement

We define a \emph{bicolor} number as a natural number $n$ with only two different digits that are repeated in two blocks (or "colors"). More formally, the sequence of digits of $n$ is $d_1d_2\ldots d_ke_1e_2\ldots e_l$, where $d$ and $e$ are the two digits and, $d \neq e$, $k > 0$ and $l > 0$.

\medskip

Examples of bicolor numbers: 7722, 44111, 666699, 277, and 45. \\
Examples of number which are \textbf{not} bicolor: 121, 113311, 7878, 1234, 7, 99910.

\medskip

Implement a \textbf{function} \verb|is_bicolor| that receives a natural number and determines if it is bicolor. The function receives a number $n > 0$ and returns \texttt{true} if it is bicolor and \texttt{false} otherwise.

\medskip

The function header must be exactly:

\begin{verbatim}
/**
* @pre n >= 0
* @post returns true if n is bicolor, false otherwise
*/
bool is_bicolor(int n);
\end{verbatim}

\Observation

You only need to submit the requested function; the main program will be ignored.
